\documentclass{USN-MSc}

\usepackage[style=apa, backend=biber]{biblatex}
\addbibresource{thesis.bib}

% --- Fyll inn forsideinformasjon ---
\faculty{Fakultet for\dots.}
\subjectinfo{Emnenavn/emnekode}
\semester{Årstall og semester}
\candidate{[Navn/kandidatnummer]}
\title{[Tittel]}
\subtitle{[Undertittel]}

\begin{document}

\USNcover

% Publisher info page: institutt, årstall, forfatter, studiepoeng
\USNpublisherinfo{Institutt for [\dots]}{[Årstall]}{[Forfatternavn]}{[antall]}

% =====================================================
\USNsection{Sammendrag}

Sett inn en oppsummering av avhandlingen.

% =====================================================
\USNsection{Abstract}

Same as in sammendrag, but in English (optional).

% =====================================================
\tableofcontents

% =====================================================
\USNsection{Forord}

Forordet skal gi plass til opplysninger og kommentarer som ikke hører hjemme i
selve avhandlingen, men som forfatteren synes er viktig å få fram. Det kan
være praktiske opplysninger («avhandlingen krever noe kjennskap til
programmering»), leserveiledning («avhandlingen har et vedleggshefte») og
omtale av personer eller institusjoner («vi vil gjerne takke firmaet A/S Fort
og Gæernt for samarbeidet»).

\bigskip
\noindent\textit{<Sted, dato>}\\
\textit{<Studentnavn>}

% =====================================================
\section{Innledning}

Dette er Universitetet i Sørøst-Norges mal for mastergradsavhandlinger. Teksten
i malen erstatter du med din egen.

Malen inneholder en del forhåndsdefinerte stiler som hjelper deg til å få et
ensartet utseende på overskrifter og på tekst gjennom hele avhandlingen.
Elementer som sidetall, innholdsfortegnelse og annet er også lagt inn slik at
du ikke skal behøve å formatere dette selv.

\subsection{Noen trinn for å lage en god mastergradsavhandling}

Det kan være lurt å ta en utskrift av malen før du begynner for å beholde
denne veiledningen når du begynner å endre på innholdet.

\subsection{Stiler}

Denne malen bruker en del stiler for å formatere (sørge for utseende på)
avsnitt. Stiler er forhåndsdefinerte oppsett for all slags utseende på teksten
i avsnitt, som skrifttype (font), skriftstørrelse, fet, kursiv, hevet, senket,
avstand mellom linjene, innrykk osv.

En av de store fordelene med å bruke stiler, er at alle avsnitt som har samme
stil, vil ha samme utseende (format) -- for eksempel overskrifter.

\subsection{Fotnoter}

Fotnoter\footnote{Slik ser en fotnote ut.} er en måte å gi tilleggsinformasjon
som ikke passer i hovedteksten. Disse nummereres automatisk.

\subsection{Litteraturhenvisninger}

Litteraturhenvisninger kan du gjøre slik: \parencite{seip1993}. Eller bare rett
i teksten \parencite{bombadil1999}. Henvisninger listes så i litteraturlisten
bakerst.

\subsection{Henvisninger til KI}

Dersom det er brukt KI i arbeidet med masteroppgave er det naturlig at det
skrives en mer omfattende tekst om hvordan KI er brukt og hvordan det har
påvirket arbeidet med oppgaven. Lag en egen del i oppgaven, for eksempel med
overskriften ``Bruk av kunstig intelligens'', eller legg det som fotnote.

Se mer informasjon om USN sine retningslinjer for bruk av KI her:
\url{https://www.usn.no/studier/studiekvalitetsuka/eksamen-fusk/retningslinjer-for-bruk-av-kunstig-intelligens-ki-ved-eksamen-og-studentoppgaver}

\subsection{Tabeller}

Under følger et eksempel på en tabell med tabelltekst
(Tabell~\ref{tab:august2005}). Tabeller skal nummereres. Bruk gjerne
\texttt{booktabs}-pakken (\texttt{\textbackslash toprule},
\texttt{\textbackslash midrule}, \texttt{\textbackslash bottomrule}) for et
profesjonelt utseende.

\begin{table}[H]
  \centering
  \caption{Kalendermåneden august 2005}
  \label{tab:august2005}
  \begin{tabular}{ccccccc}
    \toprule
    Mandag & Tirsdag & Onsdag & Torsdag & Fredag & Lørdag & Søndag \\
    \midrule
    1  & 2  & 3  & 4  & 5  & 6  & 7  \\
    8  & 9  & 10 & 11 & 12 & 13 & 14 \\
    15 & 16 & 17 & 18 & 19 & 20 & 21 \\
    22 & 23 & 24 & 25 & 26 & 27 & 28 \\
    29 & 30 & 31 &    &    &    &    \\
    \bottomrule
  \end{tabular}
\end{table}

\subsection{Figurer og bilder}

Figurer skal nummereres og ha figurtekst plassert under figuren
(Figur~\ref{fig:usnlogo}). Med \texttt{subcaption}-pakken kan du også sette
opp delfigurer.

\begin{figure}[H]
  \centering
  \includegraphics[width=0.4\textwidth]{USN_logo_main_nb_sort}
  \caption{USN-logoen brukt i denne malen.}
  \label{fig:usnlogo}
\end{figure}

\subsection{Formler}

Formler skal være innrykket og nummerert:
\begin{equation}
  (x+a)^n = \sum_{k=0}^{n} \binom{n}{k} x^k a^{n-k}
  \label{eq:binomial}
\end{equation}

For SI-enheter kan du bruke \texttt{siunitx}: \SI{9.81}{\meter\per\second\squared},
\SI{42}{\kilo\gram}.

\subsection{Kode}

Kodelistinger settes med \texttt{listings}-pakken:

\begin{lstlisting}[language=Python, caption={Eksempel på Python-kode.}, label={lst:hello}]
def hello(name):
    print(f"Hei, {name}!")

hello("verden")
\end{lstlisting}

Se Listing~\ref{lst:hello}.

\subsection{Lister}

Med \texttt{enumitem} kan lister tilpasses:
\begin{itemize}[leftmargin=*, itemsep=0.2em]
  \item Punktliste med justert venstremargin.
  \item Andre punkt.
\end{itemize}

\begin{enumerate}[label=\arabic*., leftmargin=*]
  \item Nummerert liste.
  \item Andre punkt.
\end{enumerate}

% =====================================================
\section{Om resten av malen}

Resten av malen inneholder et forslag til en disposisjon for en
mastergradsavhandling med hovedoverskrifter. Hovedoverskriftene har stilen
``Overskrift~1'' og de følges av vanlig tekst.

% =====================================================
\section{Problemstilling}

Vanlig tekst.

% =====================================================
\section{Metoder}

Vanlig tekst.

% =====================================================
\section{Resultater}

% =====================================================
\section{Diskusjon}

% =====================================================
\section{Konklusjon}

% =====================================================
\USNsection{Referanser/litteraturliste}

Alle avhandlinger skal ha en litteraturliste. Denne sorteres vanligvis
alfabetisk på forfatter.

\nocite{*}
\printbibliography[heading=none]

% =====================================================
\USNsection{Oversikt over tabeller og figurer}

\listoftables
\listoffigures

% =====================================================
% Switch numbering to "Vedlegg A", "Vedlegg B", ...
\USNappendix

\section{Vedleggstittel}

Vedlegg kan være forsøksresultater, tegninger, diagrammer, spørreskjemaer
eller lignende. Store datamengder bør plasseres i vedleggene, ellers kan
leseren gå trøtt av detaljer.

\section{Annet vedlegg}

Vedlegg som har flere sider skal sidenummereres.

\end{document}
